\documentclass[12pt]{article}
\usepackage[utf8]{inputenc}
\usepackage[T1]{fontenc}
\usepackage{hyperref}
\title{Warrgamay Language Resources for K-12 Curriculum Development}
\author{Compiled by Claude Sonnet 5.5 (claude-sonnet-5-5) for Andr\'as Kornai}
\date{}
\begin{document}
\maketitle

\section{Introduction}
Warrgamay (Wargamay, Warakamai; wgy) is a Pama-Nyungan language of northeast Queensland, closely related to Dyirbal. It is dormant, and sources describe dialects around the Herbert River, Halifax Bay and Hinchinbrook Island (including Biyay). The Endangered Languages database lists it as critically endangered, on speaker data from 1981. AIATSIS records no word list, grammar or text collection for the Warrgamay/Biyay entry.

\section{Primary Educational Resources}
The State Library of Queensland publishes a Warrgamay body chart with body-part words, and the same page points to a language pen for hearing the words. It also refers to digital stories featuring Warrgamay on ABC Open, but no direct URL for those was seen.
\begin{itemize}
\item SLQ Warrgamay body chart (PDF): \url{https://content.slq.qld.gov.au/sites/default/files/0011-211889-gamburu-gamu-biyu-warrgamay_0.pdf}
\end{itemize}

\section{Linguistic and Reference Materials}
The main scholarly description is R. M. W. Dixon, ``Wargamay'', in \emph{Handbook of Australian Languages}, vol. 2 (ANU Press, 1981). I did not see a URL for it in the searches.
\begin{itemize}
\item AIATSIS language entry Y220 Biyay: \url{https://collection.aiatsis.gov.au/node/1510719}
\item Wikipedia, Warrgamay language: \url{https://en.wikipedia.org/wiki/Warrgamay_language}
\item Endangered Languages Project entry: \url{https://www.endangeredlanguages.com/node/88838}
\end{itemize}

\section{Religious and Community Resources}
\begin{itemize}
\item Warrgamay Bada Bada Aboriginal Corporation RNTBC (native title body): \url{https://www.wettropicsplan.org.au/contacts-plan-partners/warrgamay-bada-bada-aboriginal-corporation-rntbc}
\item Girringun Aboriginal Corporation (nine tribal groups including Warrgamay): \url{https://www.wettropicsplan.org.au/contacts-plan-partners/girringun-aboriginal-corporation/}
\end{itemize}
No language programme run by either body was found.

\section{Digital and Audio Resources}
No audio archive was found. The language pen mentioned by SLQ is the only audio-related item seen.

\section{Curriculum Development Context}
The language is not known to be taught in schools. A revitalization curriculum would have to be built from Dixon's 1981 description and archival material. It should be developed with the Warrgamay community bodies above, who hold the authority over it. The resources found are enough only for a small starter unit, such as body-part vocabulary.
\end{document}